% META: {"id": "TNSI-ALGO-COURS-01", "chapitre": "TNSI-ALGORITHMIQUE", "type_objet": "cours", "capacites_codes": ["C1"], "sources_inspiration": [], "mode_creation": "ex_nihilo", "status": "approved", "capacites": ["TNSI-ALGORITHMIQUE-C1"]}
\section{Taille, hauteur et parcours d'un arbre binaire}

% ============================================================
% STRATE 1 — Taille et hauteur
% ============================================================

On représente un arbre binaire par des nœuds chaînés :
\begin{python}
class Noeud:
    def __init__(self, valeur, gauche=None, droite=None):
        self.valeur = valeur
        self.gauche = gauche
        self.droite = droite
\end{python}
Un arbre vide est représenté par \lstinline{None}.

\definition[1]{
La \textbf{taille} d'un arbre est son nombre de nœuds. La \textbf{hauteur} d'un arbre est la longueur du plus long chemin de la racine à une feuille : par convention, la hauteur de l'arbre vide est $-1$, et la hauteur d'une feuille (nœud sans enfant) est $0$.
}

\margeAppui{
Ces deux mesures se calculent par des fonctions \textbf{récursives} : le cas de base est l'arbre vide (\lstinline{None}), et le cas général combine les résultats obtenus sur les sous-arbres gauche et droit.
}

\textbf{Exemple.}
\begin{python}
def taille(arbre):
    if arbre is None:
        return 0
    return 1 + taille(arbre.gauche) + taille(arbre.droite)

def hauteur(arbre):
    if arbre is None:
        return -1
    return 1 + max(hauteur(arbre.gauche), hauteur(arbre.droite))
\end{python}

% BEGIN-VERIFY
% class Noeud:
%     def __init__(self, valeur, gauche=None, droite=None):
%         self.valeur = valeur
%         self.gauche = gauche
%         self.droite = droite
%
% def taille(arbre):
%     if arbre is None:
%         return 0
%     return 1 + taille(arbre.gauche) + taille(arbre.droite)
%
% def hauteur(arbre):
%     if arbre is None:
%         return -1
%     return 1 + max(hauteur(arbre.gauche), hauteur(arbre.droite))
%
% feuille = Noeud(5)
% assert taille(feuille) == 1
% assert hauteur(feuille) == 0
% arbre = Noeud(1, Noeud(2, Noeud(4)), Noeud(3))
% assert taille(arbre) == 4
% assert hauteur(arbre) == 2
% assert taille(None) == 0
% assert hauteur(None) == -1
% END-VERIFY

% ============================================================
% STRATE 2 — Parcours en profondeur et en largeur
% ============================================================

\propriete[Parcours en profondeur]{
Un parcours \textbf{en profondeur} visite récursivement les sous-arbres avant de \og remonter \fg{}. Trois ordres sont usuels selon la position de la racine :
\begin{itemize}
  \item \textbf{préfixe} : racine, puis sous-arbre gauche, puis sous-arbre droit ;
  \item \textbf{infixe} : sous-arbre gauche, puis racine, puis sous-arbre droit ;
  \item \textbf{suffixe} : sous-arbre gauche, puis sous-arbre droit, puis racine.
\end{itemize}
}

\textbf{Exemple.}
\begin{python}
def infixe(arbre):
    if arbre is None:
        return []
    return infixe(arbre.gauche) + [arbre.valeur] + infixe(arbre.droite)
\end{python}

% BEGIN-VERIFY
% class Noeud:
%     def __init__(self, valeur, gauche=None, droite=None):
%         self.valeur = valeur
%         self.gauche = gauche
%         self.droite = droite
%
% def infixe(arbre):
%     if arbre is None:
%         return []
%     return infixe(arbre.gauche) + [arbre.valeur] + infixe(arbre.droite)
%
% def prefixe(arbre):
%     if arbre is None:
%         return []
%     return [arbre.valeur] + prefixe(arbre.gauche) + prefixe(arbre.droite)
%
% def suffixe(arbre):
%     if arbre is None:
%         return []
%     return suffixe(arbre.gauche) + suffixe(arbre.droite) + [arbre.valeur]
%
% arbre = Noeud(1, Noeud(2, Noeud(4), Noeud(5)), Noeud(3))
% assert infixe(arbre) == [4, 2, 5, 1, 3]
% assert prefixe(arbre) == [1, 2, 4, 5, 3]
% assert suffixe(arbre) == [4, 5, 2, 3, 1]
% END-VERIFY

\margeAppui{
Un parcours \textbf{infixe} d'un arbre binaire de recherche (voir section suivante) affiche toujours les valeurs dans l'\textbf{ordre croissant} : c'est une propriété très utile pour vérifier qu'un arbre est bien un ABR.
}

\definition[2]{
Un parcours \textbf{en largeur} (\textit{Breadth-First Search}, BFS) visite les nœuds \textbf{niveau par niveau}, de la racine vers les feuilles. Il s'implémente avec une \textbf{file} : on enfile la racine, puis, tant que la file n'est pas vide, on défile un nœud, on le traite, et on enfile ses enfants non vides.
}

\begin{python}
def largeur(arbre):
    if arbre is None:
        return []
    resultat = []
    file = [arbre]
    while file:
        noeud = file.pop(0)
        resultat.append(noeud.valeur)
        if noeud.gauche is not None:
            file.append(noeud.gauche)
        if noeud.droite is not None:
            file.append(noeud.droite)
    return resultat
\end{python}

% BEGIN-VERIFY
% class Noeud:
%     def __init__(self, valeur, gauche=None, droite=None):
%         self.valeur = valeur
%         self.gauche = gauche
%         self.droite = droite
%
% def largeur(arbre):
%     if arbre is None:
%         return []
%     resultat = []
%     file = [arbre]
%     while file:
%         noeud = file.pop(0)
%         resultat.append(noeud.valeur)
%         if noeud.gauche is not None:
%             file.append(noeud.gauche)
%         if noeud.droite is not None:
%             file.append(noeud.droite)
%     return resultat
%
% arbre = Noeud(1, Noeud(2, Noeud(4), Noeud(5)), Noeud(3))
% assert largeur(arbre) == [1, 2, 3, 4, 5]
% END-VERIFY

\erreurFrequente{
\textbf{Confondre parcours en profondeur et en largeur.} Un parcours en profondeur (préfixe/infixe/suffixe) descend au maximum dans un sous-arbre avant de passer au suivant ; un parcours en largeur traite \textbf{tous} les nœuds d'un niveau avant de passer au niveau suivant. Sur l'arbre \lstinline{Noeud(1, Noeud(2, Noeud(4), Noeud(5)), Noeud(3))}, le parcours préfixe donne \lstinline{[1, 2, 4, 5, 3]} alors que le parcours en largeur donne \lstinline{[1, 2, 3, 4, 5]}.
}
